\section{Conclusions}

The model and the CCUS application together addressed one stage-gate reporting problem. Members sharing the same project information may weight adverse scenarios differently and classify the prevailing price differently. A compressed report hides whether they do. The analytical results, the benchmark evaluation, and the CCUS reconstruction jointly support one central finding. Within the maintained model, the member-threshold vector and the contested-price interval showed at each gate whether all members classified the prevailing price alike. A unique critical weight partitioned these thresholds around the no-ambiguity anchor, and trigger curvature determined its location.

Under the affine rule, the side of one half on which the critical weight fell followed from trigger curvature in the declared ambiguity coordinate. Under the pre-commitment benchmark, it also depended on the reference project value. In standard-deviation units, the affine critical weight lay above one half throughout the domain. A symmetric declaration in variance units at the same relative radius selected different endpoint volatilities and reversed this placement. Evaluating the pre-commitment benchmark at $V_0=I$ also reversed it.

Across the 33 admissible calibrations, the maximum member-trigger error was 2.74\% of the benchmark over the primary radius range up to 22\%. It reached 6.15\% with the 33\% stress radius included. Small trigger errors did not guarantee identical classifications. The two prespecified fidelity criteria held only at the evaluated radii up to 10\%. On held-out configurations, the price-axis distance used for release behaved as on the evaluated grid. The normalized criteria failed on one weight profile clustered near one half. There, a similar absolute displacement formed a much larger fraction of the narrow benchmark interval. The release rule therefore conditions on absolute distances.

In the CCUS application, baseline unanimity depended on the configuration. Under the reference configuration, none of the 1,073 archived allowance prices reached the interval in force at the time. One classification summary would therefore have sufficed at every gate. With the price record fixed and only the threshold side perturbed, this outcome failed in 15 of the 32 sensitivity configurations. Eleven of these 15 cells went to benchmark recomputation, ten of them at radii above the affine-only cap. The remaining four were binding and releasable on the affine report alone. The results have two practical implications. In this illustration, member disagreement moved the reported edges less than the common appraisal inputs did. A committee with an unsettled cost basis would weigh member attitudes inside an interval narrower than the shift from a 10\% cost revision. Repeated unanimous classifications may also arise with dispersed member thresholds, because the prevailing price can stay outside the contested interval. Unanimity alone therefore says little about the closeness of the thresholds. These implications hold under the maintained model, the reference profile, and the declared ambiguity coordinate.

The most important limitation is the maintained price model. Every member threshold is derived under a perpetual geometric Brownian motion with ambiguity on volatility alone. This assumption matters because the threshold levels, and hence the interval edges, depend on the price process. Under mean reversion, the anchor threshold moved by more than the reference margin. The report remains usable in this study for two reasons. First, the archived prices did not reject the persistence assumed by the model. Second, in the robustness grids, the critical weight stayed above one half under mean reversion, a finite horizon, and jumps at fixed declared variance. A gate expecting another process recalibrates before release, and the registry withholds affine-only release for unregistered configurations. Committee-level behavioral data could further test whether recorded weights match the positions of members before deliberation.

The proposed report therefore moves the detection of member disagreement to an earlier stage of the stage-gate process. By using one recorded weight per member and the prevailing price, it shows whether the committee agrees before any vote is taken. Where the members disagree, the report locates the disagreement in price units and states whether the affine rule can stand alone.

\section*{Declaration of competing interest}

The authors declare that they have no known competing financial interests or personal relationships that could have appeared to influence the work reported in this paper.

\section*{Data availability}

The anonymized replication archive accompanying this submission contains the processed data and the code. Its README itemizes the contents, and a per-file SHA-256 manifest lists every file. The 74 commands listed in the README were run from a clean copy. Of these, 68 completed, and 5 were skipped because their inputs are not redistributed. One legacy diagnostic returned its expected nonzero status. The code is released under the MIT License and the processed data under CC BY 4.0. The archive is a static snapshot of the evaluated registry cells and the fallback logic. The 682 policy documents are third-party material from the PKULaw database and are not redistributed. The archive lists the stable PKULaw document number of each record, so a reader with database access can retrieve the same set. One environment variable points the pipeline to a local copy. The carbon prices are published CEA and Shanghai pilot closing prices. The encoder weights are publicly released models used under their own licenses. The archive includes the processed price series and the derived document vectors. The submission-stage archive has no permanent external repository identifier yet.

\section*{Declaration of generative AI and AI-assisted technologies in the manuscript preparation process}

During the preparation of this work, the authors used a generative AI assistant for language editing, sentence restructuring, and manuscript organization. After using this tool, the authors reviewed and edited the content as needed. They independently verified the scientific claims, the derivations, the numerical results, and the references. The authors take full responsibility for the content of the article. The authors specified and executed the models, data processing, code, and analyses. The replication archive documents every command producing a reported number.
